\section{Different measurement instruments: the invisible bias}\label{sec:instruments}

A variable missing in one channel is the most visible case of a general problem: each channel measures with its own instrument. Keeping the notation of Section~\ref{sec:channel}, let $X^*$ be the latent covariate, taking values in a standard Borel space $\X^*$, with laws $P^*$ and $Q^*$ in the target group and the donor channel. Each channel records $X=\varphi_D(X^*)$, with $\varphi_0,\varphi_1\colon\X^*\to\X$ measurable, so that $P=P^*\circ\varphi_1^{-1}$ and $Q=Q^*\circ\varphi_0^{-1}$. Let $\mu^*(x^*)=\E[Y(0)\mid X^*=x^*,D=0]$.

\begin{theorem}[Different measurement instruments]\label{thm:measurement}
Assume (A1), (A4), latent support $P^*\ll Q^*$ and latent (A2): $\E[Y(0)\mid X^*,S=1]=\mu^*(X^*)$ $P^*$-almost surely.
\begin{enumerate}
  \item \emph{Common variables.} For every measurable map $\xi$ on $\X^*$, $P^*\circ\xi^{-1}\ll Q^*\circ\xi^{-1}$. In particular, if a measurable $\upsilon$ on $\X$ satisfies $\upsilon\circ\varphi_0=\upsilon\circ\varphi_1$ (variables both channels measure the same way), the laws of $U=\upsilon(X)$ in the two groups satisfy $P_U\ll Q_U$.
  \item \emph{Identification with the common variables.} If, in addition, $\mu^*(X^*)=\E[Y(0)\mid U,D=0]$ $Q^*$-almost surely, Theorems~\ref{thm:identification}--\ref{thm:lebesgue} apply with $U$ as covariates.
  \item \emph{Differential measurement with support.} If $\varphi_0$ is injective and $P\ll Q$ (on standard Borel spaces, $\varphi_0$ then has a Borel image and a measurable inverse by the Lusin--Souslin theorem \citep[Corollary 15.2]{kechris1995}, and $\varphi_1(X^*)$ lies in that image $P^*$-almost surely because $Q$ is concentrated on it) and $Y$ is bounded, the weighting functional with $X$ as covariates, $\E[w(X)\,Y\mid D=0]$ with $w=dP/dQ$, equals
  \[
  \theta_X=\int_{\X^*}\mu^*\bigl(\varphi_0^{-1}(\varphi_1(x^*))\bigr)\,P^*(dx^*),\qquad
  \theta_X-\theta=\int_{\X^*}\bigl[\mu^*\circ\varphi_0^{-1}\circ\varphi_1-\mu^*\bigr]\dd P^*.
  \]
  \item \emph{Indistinguishability.} Under the conditions of 3, there is a world without measurement error, with $\varphi_1=\varphi_0$ and the latent law of the target group equal to $P^*\circ(\varphi_0^{-1}\circ\varphi_1)^{-1}$, that reproduces the law of the target-group and donor-channel data (and of all observed data if, in addition, the law of the recorded covariates in the whole premium channel is absolutely continuous with respect to $Q$), satisfies the maintained assumptions and has $\theta=\theta_X$. No procedure based only on the target-group and donor-channel data (or, under the additional condition, on all observed data) distinguishes $\theta$ from $\theta_X$.
\end{enumerate}
\end{theorem}
\begin{proof}
(1) If $Q^*(\xi^{-1}(B))=0$, latent support gives $P^*(\xi^{-1}(B))=0$. If $\upsilon\circ\varphi_0=\upsilon\circ\varphi_1$, then $U=\upsilon(\varphi_D(X^*))=(\upsilon\circ\varphi_0)(X^*)$ in both channels, and it suffices to apply the above with $\xi=\upsilon\circ\varphi_0$.

(2) The proof of Proposition~\ref{prop:drop} uses the product structure $X^*=(V,C^*)$ only through the facts that $V$ is a measurable function of $X^*$ and is recorded identically in both channels; both hold for $U=(\upsilon\circ\varphi_0)(X^*)$, and part 1 gives the support of $U$.

(3) In the donor channel, $X=\varphi_0(X^*)$ and $\varphi_0$ is injective with a measurable inverse, so conditioning on $X$ is equivalent to conditioning on $X^*$; by (A1), $m_0(x)=\mu^*(\varphi_0^{-1}(x))$ $Q$-almost surely, and since $P\ll Q$ this version is fixed $P$-almost surely. By Steps 3 and 4 of the proof of Theorem~\ref{thm:identification}, which use only $P\ll Q$ and integrability, $\E[w(X)Y\mid D=0]=\int m_0\dd P$, and by the change of variables for image measures, $\int m_0\dd P=\int m_0(\varphi_1(x^*))\,P^*(dx^*)=\theta_X$. By latent (A2) and the tower property, $\theta=\int\mu^*\dd P^*$.

(4) In the alternative world the donor channel does not change. The target group is recorded with $\varphi_0$ and its latent law is $\tilde P^*=P^*\circ(\varphi_0^{-1}\circ\varphi_1)^{-1}$, so its observed law is $P^*\circ(\varphi_0\circ\varphi_0^{-1}\circ\varphi_1)^{-1}=P^*\circ\varphi_1^{-1}=P$, the same as in the real world. For its observed outcome it suffices to take, given $X$, the same conditional law as in the real world, and under latent (A2) its response in the basic channel is $\mu^*$; its $\theta$ is $\int\mu^*\dd\tilde P^*=\theta_X$. The alternative world satisfies the assumptions: since $\varphi_0$ is injective, $Q^*(A)=Q(\varphi_0(A))$ and $\tilde P^*(A)=P(\varphi_0(A))$, so $P\ll Q$ gives $\tilde P^*\ll Q^*$, and latent (A2) holds by construction. Under the additional condition, premium units outside the target group are treated in the same way, transforming their latent law by $\varphi_0^{-1}\circ\varphi_1$, which is then defined almost surely; without it, $\varphi_1(X^*)$ may fall outside the image of $\varphi_0$ for those units, and a world with $\varphi_1=\varphi_0$ cannot reproduce them.
\end{proof}

By part 3, $\theta_X$, not $\theta$, is the functional that weighting estimators target: one whose normalized weights converge to $w$ in $L^1(Q)$ converges to $\theta_X$ for bounded $Y$. The doubly robust estimator targets the same functional when either its weights or its outcome model is correct and the other converges to a fixed integrable limit, because the algebra of Proposition~\ref{prop:dr}, with $\int m_0\dd P$ in place of $\theta$, uses only $P\ll Q$. Part 4 is a form of the untestability of (A2), but with a specific content: the alternative world has no measurement error and the same outcome model $\mu^*$, so the bias $\theta_X-\theta$ is indistinguishable from a real difference in the latent covariates of the target group. In applications $\varphi_0$ is injective only up to the recording precision (BMI rounded to 0.1, say), which we ignore.

\begin{remark}[Stochastic instruments]\label{rem:stochastic}
Parts 3 and 4 use the target group's instrument only through the law $P$ of its recorded covariates. They therefore hold verbatim when that instrument is stochastic, $X=\varphi_1(X^*,\zeta)$ with $\zeta$ an auxiliary random variable such as the error of a self-report: if $\varphi_0$ is injective and $P\ll Q$, the weighting functional equals $\theta_X=\E[\mu^*(\varphi_0^{-1}(X))\mid S=1]$, and the world in which the latent covariates of the target group are distributed as $\varphi_0^{-1}(X)$ given $S=1$ reproduces the target-group and donor data. No condition on how $\zeta$ relates to $Y(0)$ is needed, because the limit involves the outcome only through $m_0$, which the donor channel fixes.
\end{remark}

\paragraph{Reading.} A variable measured differently is more dangerous than a missing one. The missing one breaks support, and the ideal classifier shows it by being certain, $p=1$; estimated classifiers showed it clearly in our simulation ($\hat\varepsilon>0$, Section~\ref{sec:channel}) but only faintly at the edge of a ceiling (Section~\ref{sec:lalonde}). The differently measured one can preserve support, and the classifier interprets the measurement difference as a real difference between units. In an example with $P^*=N(1,1)$, $Q^*=N(0,1)$, $\mu^*(x)=\sigma(-3-0.8x)$, with $\sigma(t)=1/(1+e^{-t})$, and a premium channel that records the variable shifted by $\Delta$, a shift of $\Delta=\pm0.3$ standard deviations produces biases of $-20.5\,\%$ and $+25.4\,\%$ on $\theta=2.92\,\%$. In a simulation of this example with 400{,}000 donor and 100{,}000 target observations and 20 replications (code and output in the repository), the weighting estimator with a logistic classifier matches $\theta_X$ ($z=+0.6$ in both cases). The population AUC, $\Phi(\delta/\sqrt2)$ for a mean difference $\delta$, changes (0.76 without error; 0.82 and 0.69 with it), but in a way that a real population difference would also produce. The only way out is information external to the design: units measured with both instruments, for instance customers who switched channels, which allow estimating $\varphi_0^{-1}\circ\varphi_1$, as in the external calibration of \citet{webbvargas2017}.

\subsection{Evidence from public data: self-reported versus measured body mass index}\label{sec:nhanes}

Theorem~\ref{thm:measurement} can be checked directly with data in which both instruments exist for the same people.

\paragraph{Data.} We use the 2017--2018 cycle of the US National Health and Nutrition Examination Survey \citep{nchs2020}: the 4{,}034 adults aged 20 or more with measured and self-reported height and weight, systolic blood pressure (mean of up to three readings), waist circumference, education and income-to-poverty ratio. We do not use the survey weights. The 4{,}034 adults form the finite population over which each target group is drawn; the theory, however, concerns a superpopulation with smooth regressions, of which they are a sample. Literally, in the finite population the recorded profiles $(V,\text{self-reported BMI})$ of the target group rarely coincide with any donor profile, so $P\ll Q$ holds only in that superpopulation, and both the estimators and our prediction reach it through smooth fitted models. Self-reported BMI is computed from the reported weight and height and rounded to 0.1, the precision of the published measured BMI, for comparability. Self-reported BMI is on average 0.81 lower than measured BMI (standard deviation 1.99), and the gap grows with BMI: from 0.00 in the lowest quartile of measured BMI to $-1.66$ in the highest, in line with the literature on self-reported anthropometry \citep{connorgorber2007}. Using it in one channel only is differential measurement in the sense of Theorem~\ref{thm:measurement}; that the error is not classical (it grows with BMI) is what makes the bias sizeable.

\paragraph{Design.} Each replication assigns every person to the target group (a synthetic premium channel) with probability $\sigma\bigl(-1+\beta_b \tilde{\mathrm{BMI}}+0.5\,\tilde{\mathrm{age}}-0.4\,\tilde{\mathrm{income}}\bigr)$, where BMI is the measured value, the tilde denotes standardization and $\beta_b=0.8$; about 30\,\% end up in the target group. The channel is synthetic, but covariates, outcomes and measurement error are real. Since the assignment depends only on the true covariates, latent exchangeability and support hold by construction, and $\theta$, the mean outcome of the target group, is observed in each replication. We compare four specifications, all with age, sex, race, education and income: (A) measured BMI in both groups; (B) self-reported BMI in the target group and measured BMI in the donors, the channel-measurement case of Theorem~\ref{thm:measurement}; (C) without BMI; (D) self-reported BMI in both groups. The estimator is the doubly robust one with five-fold cross-fitting, with a penalized logistic classifier and a ridge regression, both additive in cubic splines; there are 200 replications. The replications share the population and only the assignment varies, so we evaluate bias, not coverage.

\paragraph{Registered predictions.} Before running, we registered the predictions below (Section~\ref{sec:estimators}). We computed the limit of specification B from Theorem~\ref{thm:measurement}(3), with $\mu$, fitted on the whole sample, an estimate of $m_0(v,b)=\E[Y\mid V=v,\mathrm{BMI}=b,D=0]$, which here equals $\E[Y\mid V=v,\mathrm{BMI}=b]$ because the assignment depends only on $(V,\mathrm{BMI})$. Averaging over the assignment probabilities $q_i$, the predicted bias is $\sum_iq_i[\mu(V_i,b_i^{\mathrm{self}})-\mu(V_i,b_i)]/\sum_iq_i$, namely $-0.47$\,mmHg for blood pressure and $-2.44$\,cm for the waist. The self-report is not a function of the measured value, so this is the stochastic-instrument case of Remark~\ref{rem:stochastic}. The registered criterion was ``$|z|<3$, tolerating a model error of up to 25\,\% of the predicted bias'', which we operationalized, as the robustness files state explicitly, as $|z|<3$ or a deviation within 25\,\% of the prediction; under the alternative reading, $|z|<3$ after allowing a 25\,\% band, the LightGBM blood-pressure cell below would also pass, so our count is the more conservative. The $z$ values are conditional on the fitted $\mu$. Because the prediction uses a fitted $\mu$ of the same model class as the DR outcome model, agreement partly reflects that shared model, a point we return to below. We also predicted no bias in A, no mass of $\hat p$ close to 1 in A or B, an AUC difference between A and B below 0.02, and in C a negative bias larger in absolute value than that of B for the waist.

\begin{table}[ht]
\centering
\caption{NHANES, base design (200 replications). Mean bias of the DR and H\'ajek estimators; Monte Carlo errors of the DR bias in parentheses. No replication had any target individual with $\hat p>0.99$. The outcome standard deviations in the population are 19.7\,mmHg and 17.1\,cm.}\label{tab:nhanes}
\resizebox{\textwidth}{!}{%
\begin{tabular}{@{}lcccccc@{}}
\toprule
& \multicolumn{2}{c}{Blood pressure (mmHg)} & \multicolumn{2}{c}{Waist (cm)} & & \\
\cmidrule(lr){2-3}\cmidrule(lr){4-5}
Specification & DR & H\'ajek & DR & H\'ajek & AUC & $\neff/n_0$ \\
\midrule
A: measured in both & $-0.11$ (0.06) & $-0.26$ & $-0.18$ (0.02) & $-0.73$ & 0.743 & 0.41 \\
B: differential & $-0.48$ (0.06) & $-0.54$ & $-2.56$ (0.02) & $-2.86$ & 0.715 & 0.50 \\
\quad prediction (Theorem~\ref{thm:measurement}) & $-0.47$ & & $-2.44$ & & & \\
C: without BMI & $-1.78$ (0.05) & $-1.79$ & $-10.70$ (0.04) & $-10.70$ & 0.653 & 0.74 \\
D: self-reported in both & $-0.47$ (0.06) & $-0.60$ & $-1.32$ (0.03) & $-1.77$ & 0.737 & 0.44 \\
\bottomrule
\end{tabular}}
\end{table}

\paragraph{Results.} Table~\ref{tab:nhanes} confirms the prediction for the DR estimator approximately: $-0.48$ against $-0.47$\,mmHg ($z=-0.3$) and $-2.56$ against $-2.44$\,cm. The waist deviation, 5\,\%, is within the registered tolerance but statistically distinguishable from zero ($z=-5.7$); net of the bias of A, $-2.37$ against $-2.44$. The bias in B is not negligible: for the waist it is 0.15 standard deviations, and about a quarter of the bias from omitting BMI altogether ($-10.70$\,cm), so adjusting for a mismeasured BMI removes only three quarters of the confounding it should remove. For blood pressure, which depends weakly on BMI, it is small (0.02 standard deviations); the two outcomes thus give a strong and a weak test of the theory.

Nothing in the observed data signals the problem. No target individual has $\hat p>0.99$, and $\neff/n_0$ \emph{improves} from 0.41 to 0.50: the target group, selected for high BMI, reports a BMI shifted down toward the donors, so the groups separate less and the weights look better precisely when the estimator is biased (noise alone, in D, moves $\neff/n_0$ only to 0.44). The AUC falls from 0.743 to 0.715, but specification A is not available in practice, so this fall cannot be observed. Two registered predictions fail: A has a bias of $-0.18$\,cm in the waist ($z=-8.3$), which grows with the strength of the assignment and vanishes with LightGBM ($-0.01$; Table~\ref{tab:nhanes-robust}), so it comes from the spline models; and the AUC falls by more than the registered 0.02. The H\'ajek estimator, to which Theorem~\ref{thm:measurement}(3) literally refers, has a larger finite-sample bias already in A ($-0.73$\,cm in the waist), from the estimated weights alone, and lies further from the prediction in B ($-2.86$\,cm; $-2.13$ net of A); we registered predictions for the DR estimator only.

Specification D, registered as exploratory, suggests a practical lesson with limits: using the self-reported BMI in \emph{both} groups halves the bias of the waist relative to B, but not that of blood pressure ($-0.47$ against $-0.48$). Measuring both groups with the same imprecise instrument leaves residual confounding, but does not evaluate the outcome model at a shifted argument; which effect dominates depends on the outcome. Table~\ref{tab:nhanes-robust} shows that for the waist D beat B in all four designs, by 24 to 73\,\%, whereas for blood pressure it was worse in two of them.

\begin{table}[ht]
\centering
\caption{NHANES, robustness (200 replications each). Mean bias of the DR estimator in the waist (cm) and, in the last two columns, in blood pressure (BP). For B, the registered prediction and the $z$ of the deviation are given in parentheses.}\label{tab:nhanes-robust}
\resizebox{\textwidth}{!}{%
\begin{tabular}{@{}lcccccc@{}}
\toprule
Design & A & B & C & D & B, BP (mmHg) & D, BP (mmHg) \\
\midrule
Weak assignment ($\beta_b=0.4$) & $-0.08$ & $-2.26$ ($-2.19$; $-3.8$) & $-5.69$ & $-0.61$ & $-0.42$ ($-0.41$; $-0.2$) & $-0.27$ \\
Base ($\beta_b=0.8$), splines & $-0.18$ & $-2.56$ ($-2.44$; $-5.7$) & $-10.70$ & $-1.32$ & $-0.48$ ($-0.47$; $-0.3$) & $-0.47$ \\
Strong assignment ($\beta_b=1.2$) & $-0.46$ & $-2.89$ ($-2.61$; $-10.4$) & $-14.44$ & $-2.20$ & $-0.56$ ($-0.51$; $-0.8$) & $-0.69$ \\
Base, LightGBM & $-0.01$ & $-2.42$ ($-2.44$; $+0.9$) & $-10.38$ & $-1.23$ & $-0.20$ ($-0.47$; $+4.0$) & $-0.45$ \\
\bottomrule
\end{tabular}}
\end{table}

\paragraph{Robustness.} We varied the strength with which BMI drives the assignment and replaced the splines by LightGBM in the estimators (Table~\ref{tab:nhanes-robust}). These predictions were registered after the base results had been seen and before the three new designs were run. For each design we registered five: B within the criterion above; the bias of A in the waist below 20\,\% of that of B (a relative criterion, chosen because the base run had shown a model bias in A; under $|z|<3$, A would also fail in the weak and strong designs); no target individual with $\hat p>0.99$ in A or B; C negative and larger than B in absolute value in the waist; and D smaller than B in absolute value in the waist. Of these 15 predictions, 14 hold. In all four designs, no target individual has $\hat p>0.99$ and $\neff/n_0$ is larger in B than in A.

Across the eight design--outcome comparisons for B, the observed bias is within the registered criterion in 7 and, conditional on the fitted $\mu$, not distinguishable from the prediction at $|z|<3$ in 4. The three spline-based waist cells deviate by 3, 5 and 11\,\% of the prediction, with $z$ between $-3.8$ and $-10.4$, and the deviation grows with the strength of the assignment, as does the model bias of A; the prediction captures the size of the bias but not to within Monte Carlo error. The failure is blood pressure with LightGBM: $-0.20$ against $-0.47$ ($z=4.0$). Blood pressure depends weakly on BMI and is noisy, so LightGBM flattens that dependence and transmits only part of the measurement error, while the prediction used $\mu$ fitted with splines. A check that was not registered supports this: with $\mu$ fitted by LightGBM, the predicted bias is $-0.31$, and the observed $-0.20$ is at $z=1.7$ from it (for the waist, $-2.40$ instead of $-2.44$). The prediction is thus a property of the fitted superpopulation model as much as of the data: when the outcome depends weakly on the mismeasured covariate, it depends on the model class.
